%% Appendix C ----------------------------------------------------------------
\chapter{Data Sources}
\label{app:data-sources}
\index{data sources}

Table~\ref{tab:data-sources} lists the archives from which the application retrieves data.
Access requires a network connection; most archives are public and need no registration, the
exception being JSOC exports, which need a registered e-mail address. When you publish results
based on these data, acknowledge the providers as their data policies request.

{\small
\begin{xltabular}{\linewidth}{@{}P{0.25\linewidth} P{0.33\linewidth} L@{}}
  \caption{Archives used by the \appname{}.}\label{tab:data-sources}\\
  \toprule \thead{Data} & \thead{Provider and address} & \thead{Used by} \\ \midrule\endfirsthead
  \toprule \thead{Data} & \thead{Provider and address} & \thead{Used by} \\ \midrule\endhead
  e-CALLISTO FITS spectra & e-CALLISTO archive, FHNW, \url{https://soleil.i4ds.ch/solarradio/data/2002-20yy_Callisto/} & Downloader, Timeline, Multi-Station Comparison \\
  e-CALLISTO burst lists & Monstein catalog, \url{https://soleil.i4ds.ch/solarradio/data/BurstLists/2010-yyyy_Monstein/} & Downloader, \gui{Burst List} \\
  Learmonth spectrograph & Australian Bureau of Meteorology, Space Weather Services, \url{https://downloads.sws.bom.gov.au/wdc/wdc_spec/data/learmonth/raw/} & Learmonth browser \\
  STEREO/SWAVES level-2 & NASA Space Physics Data Facility, \url{https://spdf.gsfc.nasa.gov/pub/data/stereo/combined/swaves/level2_cdf/} & SWAVES loader \\
  GOES XRS and SGPS & NOAA National Centers for Environmental Information, \url{https://data.ngdc.noaa.gov/platforms/solar-space-observing-satellites/} & GOES viewers, GOES overlay \\
  Dst index & World Data Center for Geomagnetism, Kyoto (final, provisional or real-time), \url{https://wdc.kugi.kyoto-u.ac.jp} & Dst viewer, Geomagnetic Response \\
  Kp index & GFZ German Research Centre for Geosciences, \url{https://kp.gfz.de} & Kp viewer, Geomagnetic Response \\
  Solar wind, IMF, SYM-H and AE (OMNI 5-min) & NASA Space Physics Data Facility, \url{https://spdf.gsfc.nasa.gov/pub/data/omni/omni_cdaweb/hro_5min/} & Geomagnetic Response \\
  Recent solar wind and IMF at L1 & NOAA NCEI DSCOVR archive; NOAA SWPC real-time solar wind, \url{https://services.swpc.noaa.gov/json/rtsw/} & Geomagnetic Response (stand-in where OMNI is not yet published) \\
  CMEs, flares, shocks, storms, SEP events & NASA CCMC DONKI, \url{https://ccmc.gsfc.nasa.gov/DONKI/} & Event Context Strip, Geomagnetic Response markers, Event Story \\
  GOES flare list & NOAA SWPC, through the Heliophysics Event Knowledgebase (HEK) & Event Context Strip, Event Story \\
  SOHO/LASCO CME catalog & CDAW Data Center, NASA GSFC, \url{https://cdaw.gsfc.nasa.gov} & CME catalog, Event Context Strip, Event Story \\
  SDO/AIA and SDO/HMI & Joint Science Operations Center, \url{https://jsoc.stanford.edu}; Virtual Solar Observatory & Solar Image Analysis, SunPy Explorer, PFSS (HMI synoptic) \\
  AIA synoptic 1024² FITS & JSOC AIA synoptic archive, \url{https://jsoc1.stanford.edu/data/aia/synoptic/} & Solar Image Analysis (\gui{FITS · reduced size}) \\
  SOHO/EIT, SOHO/LASCO, STEREO/SECCHI, GOES/SUVI, PROBA2/SWAP & Virtual Solar Observatory and mission archives, through SunPy & Solar Image Analysis, SunPy Explorer, overlay layers \\
  Helioviewer JPEG2000 images & Helioviewer Project, \url{https://api.helioviewer.org} & JP2 quicklooks, Live Preview, AIA Sun in the occulter, GCS CME Fitting, Event Figure thumbnails \\
  GONG synoptic and ADAPT maps & National Solar Observatory, through SunPy & PFSS \\
  NOAA active-region labels & Heliophysics Event Knowledgebase (HEK), through SunPy & Active-region labels \\
  Software updates & GitHub releases, \url{https://github.com/SaanDev/e-Callisto_FITS_Analyzer} & Update check \\
  Website and online guide & \url{https://ecallistoanalyzer.org} & \menu{Help > Official Website} \\
  \bottomrule
\end{xltabular}}

\begin{note}
  Archive availability changes over time. The calibrated SOHO/LASCO and SOHO/EIT archives lag
  real time by many months, the HMI synoptic series is published about one Carrington rotation
  behind, OMNI solar-wind data appear a few weeks after observation, the definitive AIA
  synoptic set about a week after, and STEREO-B data end in 2014. Downloaded files are cached locally and can be
  reused without network access.
\end{note}
